\chapter{Revistas em que o trabalho poderia ser publicado}
\label{cap:conc}
\subsection*{}


%----------------------------------------------------------------------
\section{EURASIP Journal on Audio, Speech, and Music Processing}
\label{sec:porter}
\subsection*{}

Citando o próprio jornal: "O *Journal on Audio, Speech, and Music Processing* é um periódico de acesso aberto e revisado por pares que reúne pesquisadores, cientistas e engenheiros dedicados à teoria e às aplicações do processamento de diversos sinais de áudio, com foco específico no avanço tanto da ciência da comunicação por meio da fala humana quanto dos sistemas automatizados de processamento de fala e áudio. O periódico dedica-se à publicação de pesquisas originais, mas também aceita artigos de revisão e tutoriais. Os artigos abrangem aspectos tanto teóricos quanto práticos do processamento de áudio, fala e música."
A partir daqui temos como é bem direto a relação da pesquisa  com o jornal em questão

%----------------------------------------------------------------------
\section{Journal of the Audio Engineering Society}

Citando o próprio jornal: "O Journal of the Audio Engineering Society — publicação oficial da AES — é o único periódico com revisão por pares dedicado exclusivamente à tecnologia de áudio. Publicado 10 vezes ao ano, está disponível para todos os membros e assinantes da AES.
A partir daqui temos como é bem direto a relação da pesquisa  com o jornal em questão